\section{Measure bytes per frame with the record overhead included}
The moving fixture occupies 24\,795 bytes. Its two full payloads occupy 24\,576 bytes;
the six delta payloads occupy 67. Eight record headers add 128 bytes and the file
header adds 24. Compare this with eight full records: 98\,456 bytes. State whether a
reported average includes the global header and record overhead.
\Needspace{12\baselineskip}
\begin{longtable}{p{13mm}p{24mm}p{23mm}p{30mm}p{30mm}}
  \toprule Frame & Kind  & File offset & Payload bytes & Record bytes \\\midrule
  0              & Full  & 24          & 12\,288       & 12\,304      \\
  1              & Delta & 12\,328     & 11            & 27           \\
  2              & Delta & 12\,355     & 24            & 40           \\
  3              & Delta & 12\,395     & 4             & 20           \\
  4              & Full  & 12\,415     & 12\,288       & 12\,304      \\
  5              & Delta & 24\,719     & 17            & 33           \\
  6              & Delta & 24\,752     & 7             & 23           \\
  7              & Delta & 24\,775     & 4             & 20           \\\bottomrule
\end{longtable}
The all-change fixture uses one command to change every pixel. Its delta payload
needs 12\,292 bytes: one four-byte pair plus 12\,288 XOR bytes. A full payload is four
bytes smaller. Being temporally related does not guarantee that a delta representation
saves space. The still fixture gives the opposite boundary: one full frame followed
by three four-byte skip-only payloads.

Separately measure instructions used to compute deltas, assemble commands, calculate
CRCs, validate the file and seek to a target. A vector kernel can reduce delta work
without changing the encoded size. Choosing more full frames can reduce replay work
while increasing file size. Retain both observations rather than combining them into
an unexplained claim that the codec is faster.
